\chapter{Experimental Evaluation}
\label{chap:evaluation}

\section{Evaluation Protocol}

The experimental evaluation tests the ability of the selected model to forecast the next daily three-dimensional ocean state. Each inference sample contains three consecutive context days and one target one day after the final context state. All parameters are frozen during evaluation. The same preprocessing code, training-derived standardization statistics, variable order, and target masks used during training are restored before the network is called.

The best v3 checkpoint, selected at epoch 94 through validation RMSE, is evaluated on the complete 1998 validation year and the previously unseen 1999 test year. Because the model consumes three context days, the last two days of the preceding year are prepended to each evaluation period as inputs only. The first forecast uses 30--31 December and 1 January to predict 2 January; the last uses 28--30 December to predict 31 December. Thus each non-leap evaluation year contains 364 targets, all belonging to the year being evaluated. The chronological split itself was defined in Chapter~\ref{chap:architecture}; no target crosses a split boundary.

With the repository in the \texttt{Thesis/thesis\_repo} directory and the dataset, mask, normalization statistics, and v3 checkpoint in the parent \texttt{Thesis} directory, the aggregate evaluation can be reproduced with:

\begin{lstlisting}[caption={Inference commands for the selected checkpoint.}]
python src/train.py --data-path ../glorys12_med_test_1994_2003.nc \
  --mask-path ../land_sea_mask.nc \
  --stats-path ../normalization_stats.nc \
  --reuse-stats \
  --checkpoint-path ../best_forecaster_v3.pt \
  --evaluate-validation --evaluate-test
\end{lstlisting}

The command executes forward passes over both evaluation splits; the appearance of a validation progress bar does not indicate that training has restarted. No optimizer is created for these evaluation branches and no gradient or parameter update is performed.

Physical-unit evaluation uses the same checkpoint for temperature, salinity, zonal velocity, and meridional velocity. The routine automatically chooses the closest level to the requested depth and reports both that slice and an aggregate over all depths. The same run also writes machine-readable JSON and CSV results and the annual error maps:

\begin{lstlisting}[caption={Annual physical-unit multivariable evaluation.}]
python src/train.py --data-path ../glorys12_med_test_1994_2003.nc \
  --mask-path ../land_sea_mask.nc \
  --stats-path ../normalization_stats.nc \
  --reuse-stats \
  --checkpoint-path ../best_forecaster_v3.pt \
  --evaluation-depth 0.5 \
  --evaluate-physical-validation \
  --evaluate-physical-test \
  --plot-annual-errors-test
\end{lstlisting}

The checkpoint is first loaded on the CPU so its configuration can define context length, channel dimensions, and internal normalization. The model is then constructed, weights are restored, and the complete object is moved to the selected evaluation device. This order prevents an inference command from silently evaluating a checkpoint with a model configuration taken from unrelated command-line defaults.

\section{Persistence Reference}

The neural forecast is compared with persistence. Given the latest input state $X_t$, persistence predicts

\[
\widehat{X}^{\mathrm{pers}}_{t+1}=X_t.
\]

This baseline has no trainable parameters and is evaluated over exactly the same targets, variables, depths, and masks as the neural model. It is particularly demanding at a 24-hour horizon because much of the ocean changes gradually. A neural model predicting the full state can achieve a low absolute error while still losing to persistence if it slightly smooths the input or introduces a small bias in regions that change little.

Persistence is more informative here than a constant climatology because it uses the most recent state available to the network. It also provides a direct test of whether the neural mapping has learned the daily increment rather than only the typical absolute state.

\section{Metrics and Aggregation}

Deterministic accuracy is measured from the predicted mean. For $N$ valid elements,

\begin{align}
\mathrm{RMSE} &= \sqrt{\frac{1}{N}\sum_{i=1}^{N}(y_i-\mu_i)^2},\\
\mathrm{MAE} &= \frac{1}{N}\sum_{i=1}^{N}|y_i-\mu_i|,\\
\mathrm{Bias} &= \frac{1}{N}\sum_{i=1}^{N}(\mu_i-y_i).
\end{align}

The normalized aggregate RMSE and MAE pool all four variables, all depths, all horizontal ocean points, and all dates of a split. They are dimensionless because every channel and depth has been standardized. Since such pooling can hide large differences among variables, residuals and predicted standard deviations are also converted back to the physical unit of each channel using the training standard deviation at the corresponding depth.

Skill relative to persistence is

\[
S_{\mathrm{RMSE}}=1-
\frac{\mathrm{RMSE}_{\mathrm{model}}^2}
{\mathrm{RMSE}_{\mathrm{pers}}^2}.
\]

A value above zero means that the neural mean has lower MSE than persistence; a value below zero means that persistence is more accurate. Because the score uses squared RMSE, even a modest difference between the two errors can produce a visibly negative value.

The probabilistic score is the masked Gaussian NLL defined in Chapter~\ref{chap:architecture}. It is not repeated algebraically here. Two marginal interval coverages complement it:

\begin{align}
C_{68} &= \frac{1}{N}\sum_i
\mathbf{1}(|y_i-\mu_i|\leq\sigma_i),\\
C_{95} &= \frac{1}{N}\sum_i
\mathbf{1}(|y_i-\mu_i|\leq1.96\sigma_i).
\end{align}

For a calibrated Gaussian forecast these fractions would approach 0.68 and 0.95 over a large homogeneous sample. Coverage alone is incomplete: a very wide interval can obtain high coverage without being sharp, and aggregate coverage can conceal miscalibration by location, variable, or depth. The mean predicted standard deviation is therefore reported beside coverage.

The implementation accumulates sums and counts over batches before computing final metrics. RMSE is obtained from the total squared error and total valid-point count rather than by averaging daily RMSE values. This makes every valid tensor element contribute equally, independently of which batch contains it. The consequence is that dates and shallow/deep regions with different numbers of wet cells do not receive equal group weight. No latitude-area or vertical-thickness weighting is applied, so the results describe the discrete machine-learning grid and should not be interpreted as volume-weighted heat or momentum errors.

Daily errors are temporally correlated, and the hundreds of millions of grid-point residuals are spatially correlated. The raw valid-point count therefore cannot be treated as a statistical sample size for a confidence interval. This thesis compares complete, identical splits deterministically and does not report formal significance tests. A future uncertainty analysis could use block bootstrap sampling over dates or physically defined regions to preserve part of that dependence.

\section{Annual Variable-wise Results}

The deterministic and probabilistic tables are generated directly from the JSON artifacts emitted by the evaluator. They contain temperature, salinity, $u$, and $v$ results for all depths and for the available level nearest to 0.5~m, namely 0.506~m. This avoids transcribing values from the older 362-pair evaluation. The generated file is accepted only when both inputs declare exactly 364 forecasts.

\begin{table}[p]
\centering\scriptsize
\setlength{\tabcolsep}{4pt}
\caption{Physical-unit deterministic results over 364 annual forecasts. Bias is forecast minus target.}
\label{tab:physical-deterministic-results}
\begin{tabular}{lllrrrrrr}
\toprule
Split & Variable & Depth & RMSE & MAE & Bias & Pers. RMSE & Pers. MAE & Skill \\
\midrule
Validation & Temperature ($^\circ$C) & All & 0.17763 & 0.08958 & -0.03890 & 0.09589 & 0.04502 & -2.43149 \\
Validation & Temperature ($^\circ$C) & 0.506 m & 0.41020 & 0.26862 & -0.01929 & 0.20788 & 0.14340 & -2.89367 \\
Validation & Salinity ($10^{-3}$) & All & 0.07160 & 0.02217 & 0.01211 & 0.04748 & 0.00772 & -1.27423 \\
Validation & Salinity ($10^{-3}$) & 0.506 m & 0.18973 & 0.07961 & 0.03986 & 0.08609 & 0.01908 & -3.85704 \\
Validation & Zonal $u$ (m\,s$^{-1}$) & All & 0.01752 & 0.00928 & -0.00047 & 0.01728 & 0.00882 & -0.02881 \\
Validation & Zonal $u$ (m\,s$^{-1}$) & 0.506 m & 0.03732 & 0.02631 & -0.00093 & 0.03321 & 0.02271 & -0.26288 \\
Validation & Meridional $v$ (m\,s$^{-1}$) & All & 0.01598 & 0.00856 & -0.00037 & 0.01658 & 0.00835 & 0.07097 \\
Validation & Meridional $v$ (m\,s$^{-1}$) & 0.506 m & 0.03553 & 0.02583 & 0.00032 & 0.03422 & 0.02329 & -0.07842 \\
Test & Temperature ($^\circ$C) & All & 0.18794 & 0.09455 & -0.03735 & 0.10397 & 0.04785 & -2.26732 \\
Test & Temperature ($^\circ$C) & 0.506 m & 0.41860 & 0.27683 & -0.00928 & 0.20777 & 0.14564 & -3.05908 \\
Test & Salinity ($10^{-3}$) & All & 0.07213 & 0.02271 & 0.01191 & 0.04477 & 0.00783 & -1.59633 \\
Test & Salinity ($10^{-3}$) & 0.506 m & 0.19219 & 0.07889 & 0.03923 & 0.08541 & 0.01896 & -4.06285 \\
Test & Zonal $u$ (m\,s$^{-1}$) & All & 0.01872 & 0.00980 & -0.00068 & 0.01875 & 0.00937 & 0.00313 \\
Test & Zonal $u$ (m\,s$^{-1}$) & 0.506 m & 0.03818 & 0.02686 & -0.00093 & 0.03394 & 0.02334 & -0.26529 \\
Test & Meridional $v$ (m\,s$^{-1}$) & All & 0.01700 & 0.00913 & -0.00060 & 0.01750 & 0.00889 & 0.05670 \\
Test & Meridional $v$ (m\,s$^{-1}$) & 0.506 m & 0.03645 & 0.02637 & -0.00073 & 0.03439 & 0.02382 & -0.12324 \\
\bottomrule
\end{tabular}
\end{table}

\begin{table}[p]
\centering\scriptsize
\caption{Physical-unit probabilistic diagnostics over the same annual forecasts.}
\label{tab:physical-probabilistic-results}
\begin{tabular}{lllrrr}
\toprule
Split & Variable & Depth & Mean $\sigma$ & 68\% coverage & 95\% coverage \\
\midrule
Validation & Temperature ($^\circ$C) & All & 0.122832 & 0.749025 & 0.959739 \\
Validation & Temperature ($^\circ$C) & 0.506 m & 0.445210 & 0.833820 & 0.978935 \\
Validation & Salinity ($10^{-3}$) & All & 0.025602 & 0.687881 & 0.944692 \\
Validation & Salinity ($10^{-3}$) & 0.506 m & 0.090442 & 0.746352 & 0.956752 \\
Validation & Zonal $u$ (m\,s$^{-1}$) & All & 0.013858 & 0.789721 & 0.956844 \\
Validation & Zonal $u$ (m\,s$^{-1}$) & 0.506 m & 0.037914 & 0.755272 & 0.958670 \\
Validation & Meridional $v$ (m\,s$^{-1}$) & All & 0.013154 & 0.795684 & 0.958516 \\
Validation & Meridional $v$ (m\,s$^{-1}$) & 0.506 m & 0.033937 & 0.709480 & 0.940354 \\
Test & Temperature ($^\circ$C) & All & 0.124849 & 0.738794 & 0.949917 \\
Test & Temperature ($^\circ$C) & 0.506 m & 0.449665 & 0.825928 & 0.979179 \\
Test & Salinity ($10^{-3}$) & All & 0.025908 & 0.696942 & 0.937503 \\
Test & Salinity ($10^{-3}$) & 0.506 m & 0.090543 & 0.755580 & 0.957487 \\
Test & Zonal $u$ (m\,s$^{-1}$) & All & 0.014053 & 0.774239 & 0.947604 \\
Test & Zonal $u$ (m\,s$^{-1}$) & 0.506 m & 0.038108 & 0.748713 & 0.955055 \\
Test & Meridional $v$ (m\,s$^{-1}$) & All & 0.013359 & 0.778566 & 0.948389 \\
Test & Meridional $v$ (m\,s$^{-1}$) & 0.506 m & 0.034107 & 0.705973 & 0.935753 \\
\bottomrule
\end{tabular}
\end{table}


For the 1999 test year at 0.506~m, the model's temperature MAE is 0.277~$^{\circ}$C versus 0.146~$^{\circ}$C for persistence. The salinity MAE is 0.0789 versus 0.0190 in the dataset's $10^{-3}$ unit. The current components have smaller but still positive surface MAE differences: approximately 15\% for $u$ and 11\% for $v$. Across all depths, the model slightly improves RMSE for $u$ and $v$ even though its MAE remains higher. The assessment therefore depends on both depth and error measure, as the complete table makes explicit.

\section{Annual Spatial Diagnostics}

For a variable $q$ and surface grid point $p$, let $\mathcal{T}_q(p)$ be the dates with a valid target, $N_q(p)=|\mathcal{T}_q(p)|$, and $a_{q,t}(p)=|\widehat y_{q,t}(p)-y_{q,t}(p)|$. The four annual diagnostics are
\begin{align}
E^{\mathrm{model}}_q(p) &= \frac{1}{N_q(p)}\sum_{t\in\mathcal{T}_q(p)} a_{q,t}(p),\\
s_{a,q}(p) &= \sqrt{\frac{1}{N_q(p)}\sum_{t\in\mathcal{T}_q(p)}
\left(a_{q,t}(p)-E^{\mathrm{model}}_q(p)\right)^2},\\
E^{\mathrm{pers}}_q(p) &= \frac{1}{N_q(p)}\sum_{t\in\mathcal{T}_q(p)}
\left|x_{q,t}(p)-y_{q,t}(p)\right|,\\
D_q(p) &= E^{\mathrm{model}}_q(p)-E^{\mathrm{pers}}_q(p).
\end{align}
Here $x_{q,t}(p)$ is the latest input state. The absolute value is applied to each forecast error before averaging; it is not applied again to $D_q(p)$. Thus a positive difference means a larger model MAE and a negative one means a smaller model MAE. This signed difference between MAEs is distinct from forecast bias, which averages the signed prediction-minus-target error. The population standard deviation $s_{a,q}(p)$ measures day-to-day variation in the model's realized absolute error, not its predicted Gaussian uncertainty. Land and cells without valid targets remain masked.

\begin{figure}[p]
\centering
\includegraphics[width=\textwidth]{annual-mae-maps-1999.png}
\caption{Model MAE over 364 one-day forecasts from 2 January to 31 December 1999 at 0.506~m. Each variable has its own physical-unit colour scale, capped at the 95th percentile of valid cells. Values above the cap are saturated.}
\label{fig:annual-mae}
\end{figure}

\begin{figure}[p]
\centering
\includegraphics[width=\textwidth]{annual-error-std-maps-1999.png}
\caption{Temporal standard deviation $s_{a,q}(p)$ of the model's absolute daily errors over the 364 forecasts. This measures variability in realized error magnitude and is distinct from the network's predicted uncertainty. Each variable has its own physical-unit scale capped at the 95th percentile.}
\label{fig:annual-error-std}
\end{figure}

\begin{figure}[p]
\centering
\includegraphics[width=\textwidth]{annual-persistence-mae-maps-1999.png}
\caption{Persistence MAE over the same dates, cells, variables, and depth as Figure~\ref{fig:annual-mae}. Each panel is capped at its own 95th percentile. Consequently, equal colours in the model and persistence figures do not represent equal numerical errors.}
\label{fig:annual-persistence-mae}
\end{figure}

\begin{figure}[p]
\centering
\includegraphics[width=\textwidth]{annual-mae-difference-1999.png}
\caption{Signed annual difference $D_q(p)$ between model MAE and persistence MAE. Red denotes a larger model error; blue denotes a smaller one. No absolute value is applied to the subtraction. The diverging scale is centred on zero and capped symmetrically at the 95th percentile of $|D_q|$ for each variable.}
\label{fig:annual-mae-difference}
\end{figure}

The difference map in Figure~\ref{fig:annual-mae-difference} provides the direct spatial comparison that the separately scaled MAE figures cannot. Of the 4,330 valid surface cells, $D_q(p)>0$ in 99.9\% for temperature and 100\% for salinity. The corresponding fractions are 62.9\% for $u$ and 59.9\% for $v$: the current maps contain sizeable regions where the model improves upon persistence, despite a positive domain-wide MAE difference. The strongest positive patches often lie near the wet-domain boundary, but the maps alone do not identify the underlying physical or numerical cause.

Figure~\ref{fig:annual-error-std} adds a temporal dimension to the annual mean: cells with similar MAE can differ in how consistently that error magnitude occurs. High $s_{a,q}(p)$ indicates more variable absolute daily errors, without revealing whether forecasts tend to be too high or too low. The four diagnostics therefore distinguish average magnitude, temporal variability, and relative performance against persistence; they do not constitute a formal significance test or establish the causes of the spatial patterns.

\clearpage
\subsection{Environmental relevance of accurate forecasts}

The following application perspective is a reworking of written observations provided by Cataldo Licchelli, marine biologist. An accurate forecast of potential temperature, salinity, and both horizontal velocity components would be valuable because these fields jointly describe the physical conditions experienced by marine organisms and the movement of water through their habitats. Temperature and salinity influence density and stratification; $u$ and $v$ together specify the direction and speed of horizontal flow. Forecasting only one of these fields well would not necessarily describe the coupled marine environment well.

Reliable temperature forecasts could help identify where and when marine organisms may face unusual thermal conditions. The observations highlight prolonged warm episodes as a concern for sensitive habitats, including benthic communities and seagrass meadows; warmer water can also hold less dissolved oxygen. Surface temperature contributes to exchanges of heat and moisture with the atmosphere. The environmental value of predicting it accurately is therefore not limited to obtaining a low average error: it includes representing the timing, duration, and location of thermal anomalies. A temperature forecast does not, by itself, predict mortality, species movement, or the evolution of a storm.

Salinity matters both to the physical structure of the water column and to organisms adapted to particular salinity ranges. Changes in temperature and salinity can modify stratification and hence the vertical exchange of oxygen and nutrients; marked local salinity changes may also impose osmotic stress in coastal habitats. A reliable salinity forecast could therefore help interpret changes in water masses and identify locations that merit biological monitoring. This is an argument for monitoring, not a claim that a given forecast error has caused ecological damage.

The two velocity components need to be accurate together because their combination determines the horizontal path of moving water. Currents can transport plankton, eggs, larvae, nutrients, and floating pollutants; circulation also contributes to the distribution of resources and to the connectivity of marine populations. Accurate current forecasts could support the study of dispersal, the placement of monitoring efforts, and responses to pollution at sea. In coastal settings, circulation can contribute to upwelling and the redistribution of nutrient-rich water. None of these outcomes follows from the MAE of $u$ or $v$ considered separately: transport and ecological consequences require their combined evolution to be assessed over time.

Taken together, these observations explain why a skillful four-variable forecast could be useful for environmental protection, fisheries and aquaculture planning, and the interpretation of marine ecosystem change. The forecast itself does not alter the environment; its potential impact lies in better-informed observation and decisions. The present experiment does not demonstrate those benefits. In particular, persistence remains more accurate at the selected surface depth for all four variables, and no organism, pollutant trajectory, oxygen field, or ecological response was evaluated. Application-specific validation would be needed before relying on this model for such purposes.
